\section{Limitations}
\label{sec:limitations}

We acknowledge several limitations, in roughly decreasing order of how much they would change the conclusion:

\begin{enumerate}
    \item \textbf{Single seed.} All numbers are from one training run. A multi-seed sweep would tighten the confidence interval on $\Delta\mathrm{AP}_S$ and confirm whether $+0.99$ is reproducible.
    \item \textbf{Floor and mask not isolated.} v2 changes both fixes simultaneously, so we cannot say from this experiment whether one alone is sufficient. A $2{\times}2$ ablation is the obvious next step.
    \item \textbf{Below the published RT-DETR baseline.} v2 reaches $\mathrm{AP}_S = 34.25$ vs the published $34.7$. Our claim is the $+0.99$ causal contribution, not a new state of the art --- the absolute deficit is explained by the schedule (13-epoch finetune vs the paper's 6$\times$ schedule from scratch).
    \item \textbf{Latency cost is dominated by P2.} The ${\approx}2{\times}$ slowdown over the RT-DETR baseline ($53.3 \to 25.9$ FPS) is essentially entirely from adding the stride-4 feature level. The feedback module itself, especially with the level mask, is a small fraction of inference cost. We see three concrete paths to recover the speed without losing the small-object gain: (i) drop P2 at inference and test whether the $+0.99$ contribution transfers to the original P3--P5 setup --- if it does, the baseline ${\approx}53$\,FPS is recovered for free; (ii) replace the full P2 input-projection block with a lightweight variant (halved channels or a single $1{\times}1$ projection); (iii) distill v2 into a student network restricted to P3--P5, using the full v2 as teacher and matching either the decoder logits or the refined memory. We consider (i) the highest-leverage next experiment: it tests whether P2 is doing the work or whether the feedback path alone is.
    \item \textbf{COCO only.} Domains where small objects dominate (aerial imagery, traffic surveillance) might show a larger or smaller gain; we did not test there.
\end{enumerate}
